\section{Conclusion}
\label{sec:conclusion}

Survey questions about generative artificial intelligence use and dependence modestly improved estimates of current burnout scores in this sample of 97 university students. Adding the three questions to the study-condition inputs reduced mean absolute error from 0.6183 to 0.5883 points on the 1-to-5 scale in the matched comparison. The reduction applies to the questions as a set. Average error remained higher in the High group at 0.9329 points.

For the warning task, workload-only logistic regression detected 16.0 of the 17 High cases on average. The eight-input random forest detected 14.8. Both classifiers averaged 28.0 false warnings among the 80 students whose measured score was below 4. In a separate comparison against chance, the selected forest's balanced accuracy was higher than in any of the 199 evaluations with shuffled labels (permutation \(p=0.005\)).

The warning counts are averages from five cross-validation repetitions on a single convenience sample. An independent evaluation should fix the prediction procedures in advance and test them on a new student cohort with measured BAT-S scores.
